\documentclass[../main.tex]{subfiles}  % 使用主文件的文档类

\begin{document}

\section{Introduction}

% In cryptographic systems, entropy is a fundamental concept used to quantify the randomness and uncertainty inherent in the system.
% Specifically, entropy pertains to the size and complexity of the key space, which is directly correlated with the security of the cryptographic system.
% A key with high entropy demotes a larger key space, thereby augmenting the difficulty for an attacker to recover the key through brute-force methods.
% Ideally keys in cryptographic protocols are considered ``perfect'', meaning that an attacker cannot recover the key from any ciphertext. 
% In the formal analysis of security protocols using symbolic models, such \emph{perfect encryption assumption} is commonly employed.

% However, in practice, achieving an absolutely perfect key is unattainable. 
% An attacker with adequate computational resources can potentially brute force any key from the corresponding ciphertext. 
% The recent advancements in chip technology have greatly enhanced attackers' computational capabilities, which in turn challenges the security of traditional key lengths previously deemed secure. 
% Unfortunately, some vulnerabilities in protocols associated with low-entropy keys have been exploited by attackers, exemplified by the PMKID \cite{pmkidattack} and KNOB \cite{antonioli2019knob} attacks. 
% These attacks compromise the security properties of the protocols by recovering low-entropy keys.

% There are several reasons for the adoption of low-entropy keys in modern cryptographic design. 
% First, some protocols might regard certain key lengths as secure; however, these can become susceptible to attacks over the years.
% Second, current versions of such protocols might maintain compatibility with these obsolete designs for the sake of backward compatibility.
% Third, some protocols opt for shorter keys to cater to devices with limited computational or power resources \cite{bluetooth6}. 
% Lastly, in pursuit of user-friendliness, certain protocols permit the use of user-selected passwords as keys \cite{80211i}, which may result in the choice of easily predictable password combinations.

Entropy is a fundamental concept that is directly correlated with the security of the cryptographic system.
A key with high entropy indicates a high level of difficulty for an attacker attempting to recover the key from the ciphertext.
However, in practice, modern cryptographic protocols, such as WPA2-PSK, Bluetooth Low Energy~(BLE), and Bluetooth Classic~(BC), usually tolerate low-entropy keys to accommodate user-friendly designs, maintain backward compatibility, or support resource-constrained devices.
% This introduces a security issue about recovering low-entropy keys in these protocols, such as PMKID~\cite{pmkidattack} and KNOB~\cite{antonioli2019knob} attacks.
This introduces a security issue concerning the recovery of low-entropy keys in these protocols, as demonstrated by attacks like PMKID~\cite{pmkidattack} and KNOB~\cite{antonioli2019knob}.


% Entropy is a fundamental concept directly correlated with the security of cryptographic systems, where a high-entropy key typically makes it challenging for an attacker to deduce the key from ciphertext.
% However, in practice, modern cryptographic protocols often tolerate low-entropy keys, as seen in widely-used standards such as WPA2-PSK, Bluetooth Low Energy, and Bluetooth Classic.
% This tolerance introduces potential vulnerabilities, allowing attackers to exploit lower-entropy keys in these protocols, as demonstrated by security issues like the PMKID attack on WPA2-PSK~\cite{pmkidattack} and the KNOB attack on Bluetooth~\cite{antonioli2019knob}.

% 使用符号模型分析协议，pioneered by Dolve et al.~\cite{}，是主流的确保协议设计安全的方法。
% 然而符号模型假设xxx，不能直接用来推到密钥熵。
% 现有研究xxx，他们有xxx缺陷。
% 或者，符号模型下分析协议密钥熵具有xxx挑战。
% Using symbolic models to formally analyze protocols has become a mainstream research approach, and a great number of  works\cite{basin2021emv,basin20185G,cremers2017tls13} successfully confirmed this.
The use of symbolic models to formally analyze protocol vulnerabilities, pioneered by Dolev et al.~\cite{dolev1983security}, has become a mainstream research approach.
% Many symbolic model tools, such as~\cite{armando2005avispa,blanchet2018proverif,escobar2007maude,cremers2008scyther,meier2013tamarin,cspa}, have been developed to automatically analyze protocols and discover attacks.
Numerous tools have been developed to automatically analyze protocols and identify potential attacks, including AVISPA\cite{armando2005avispa}, ProVerif~\cite{blanchet2018proverif}, Maude-NPA~\cite{escobar2007maude}, Scyther~\cite{cremers2008scyther}, \Tamarin{}~\cite{meier2013tamarin}, and others.
In symbolic models, messages and keys are abstracted as terms in term algebra, rather than as concrete bit sequences as used in computational models. 
Consequently, the entropy of keys is lost during this abstraction process, making it impossible for existing symbolic model tools to distinguish low-entropy keys or analyze the brute-force capabilities of attackers.

This raises a critical question: How to capture the brute-force behaviors of attackers within symbolic models?
Some studies\cite{shi2023blesc, claverie2023tamarin} have proposed solutions by manually identifying keys that are vulnerable to compromise and defining specific rules that allow these keys to be directly revealed to attackers.
However, these studies do not provide a general or automated method for analyzing the impact of low-entropy keys on the entire protocol. 
% For one thing, this approach requires manual analysis of the vulnerabilities inherent in the protocol keys instead of automated analysis.
% For another, manual analysis is inherently constrained, as it can only provide attacker with limited ways to compromise low-entropy keys and fails to provide a comprehensive strategy for attackers to exploit low-entropy keys effectively.
Firstly, this method relies on the manual identification of key vulnerabilities, limiting the scalability and objectivity of the analysis.
Secondly, manual analysis is inherently constrained, as it only offers attacks a limited set of predefined ways to compromise low-entropy keys.
As a result, it fails to provide a comprehensive strategy for effectively exploiting such vulnerabilities.

To automatically analyze protocols tolerating low-entropy keys, we encounter three main challenges:
% First, the symbolic analysis tools only remain the logical structure of messages or keys, and the entropy of keys is hard to be expressed and analyzed. 
% Second, brute-force attacks require three prerequisites: known encryption algorithm, known ciphertext, and no additional values~(salt value, IV, etc.) required for calculating the ciphertext. 
% These prerequisites are ambiguous, and it is hard to capture different numbers of additional values  and cryptographic operators in a single rule.
% Third, adversaries' brute force ability needs to be improved, no longer targeting a specific key, but can be freely combined with other capabilities and applied to the attack.
\begin{enumerate}[leftmargin=*]
    \item \textbf{Entropy Representation:}
    Symbolic analysis tools typically preserve only the logical structure of messages or keys, making it impossible to deduce and analyze the entropy of keys.
    Additional mechanisms are needed to describe a key's entropy in order to cover general brute-force attacks within the symbolic model.
    \item \textbf{Prerequisites for Brute-Force Attacks:}
    Brute-force attacks involve an attacker repeatedly attempting keys, relying on prior knowledge of the algorithm, its outputs, and other relevant inputs beside targeted keys~(e.g., salt values, initialization vectors). 
    Accurately and generally capturing these prerequisite conditions within a symbolic model presents a significant challenge, as these conditions vary across different protocols and keys.
    \item \textbf{Enhanced Adversarial Capabilities:}
    In existing works, an attacker is typically limited to compromising specific keys. 
    However, an attacker can attempt to brute-force any key as long as its entropy doesn't exceed its computational capacity. 
    The attacker can also deconstruct complex keys into simpler sub-keys and brute-force them individually. 
    A challenge lies in modeling this flexible brute-forcing capability within the symbolic model.
\end{enumerate}

This paper aims to bridge the existing gap by introducing a general and automated method for modeling low-entropy keys and the adversarial brute-force capabilities within symbolic analysis frameworks. 
We focus on enhancing \Tamarin{} to describe general brute-force attacks conducted by adversaries. 
\Tamarin{} is a state-of-the-art protocol verification tool widely used for the formal analysis of many protocols, including TLS\cite{cremers2017tls13}, 5G-AKA\cite{basin20185G, wang20215G}, EMV\cite{basin2021emv,radu2022emv}, etc.
We have designed a new syntax, which is a superset of the syntax of \Tamarin{}, allowing for the labeling of low-entropy variables in symbolic models. 
Additionally, we have developed an automated protocol analysis tool supporting such low-entropy annotated models, which we call EntropyVerif\cite{toolrepo}. 
This tool supports an adversary model that extends beyond the Dolev-Yao model, capable of capturing all brute-force behaviors against security protocols, including complex attacks where an attacker synthesizes and utilizes various attack capabilities.
Our approach can be used to identify security issues arising from low-entropy keys in protocols and to verify the security of protocols that tolerate low-entropy keys.
% shi2023blesc,claverie2023tamarin, jangid2023bluepe
% 有什么效果，结论，发现。

\noindent\textbf{Contributions.}
Our primary contributions are as follows:
\begin{itemize}[leftmargin=*]
    \item We relax the \emph{perfect encryption assumption} in the symbolic model. 
    Our methodology's adversary model surpasses the Dolev-Yao adversary by allowing the attacker to conduct brute-force attacks on all low-entropy keys.
    % In addition to the ability to eavesdrop on, block, modify, inject, and replay messages transmitted over the network, the adversary can conduct brute-force attacks on all low-entropy keys.
    % \item We applied our methodology to develop a protocol analysis tool that accepts models with low-entropy labels. 
    % This tool can automatically analyze vulnerabilities associated with low-entropy keys in protocols and generally capture brute-force attacks.
    \item We developed EntropyVerif, an extension of \Tamarin{}, to introduce low-entropy keys to protocol models. 
    Our tool allows us to discover vulnerabilities associated with low-entropy keys and generally capture brute-force attacks.
    \item By using our tool to analyze real-world protocols, we identified several known vulnerabilities related to low-entropy keys, including the PMKID attack and KNOB attacks.
    Our tool not only successfully rediscovered these existing vulnerabilities but also uncovered a novel attack, the unilateral multi-downgrade(UMD) attack.
    % We demonstrate its effectiveness by rediscovering these known vulnerabilities, and a novel attack, the unilateral multi-downgrade(UMD) attack. 
    \item The UMD attack significantly weakens the strength of the session keys in BC devices. 
    Even when the longest keys are used, an attacker can repeatedly downgrade one side of the two devices to obtain different segments of the session key and ultimately reconstruct the complete key.
    This attack compromises the session key with a computational complexity that can be as low as $\frac{1}{2^{116}}$ of that required for a brute-force attack.
\end{itemize}

\noindent\textbf{Overview.} 
% Section \ref{sec:2background} reviews protocols that tolerate low-entropy keys and the protocol analysis tool \Tamarin{}.
Section \ref{sec:2background} reviews the definition of entropy and brute-force attacks, protocols that tolerate low-entropy keys, and the protocol analysis tool \Tamarin{}.
Section \ref{sec:3symbolic-modeling} introduces our modeling approach for entropy-based security, detailing the key-breaking oracles. 
Section \ref{sec:4extension-and-compilation} discusses EntropyVerif, our extension of \Tamarin{}.
Section \ref{sec:5case-study} presents case studies where we apply our method to analyze real-world protocols. 
Finally, in Section \ref{sec:6results}, we present the results of our case studies, where we discover known vulnerabilities and a novel attack. 
We additionally  offer a review of related work in Section \ref{sec:7related-work} and provide a conclusion and discussion of our work in Section \ref{sec:8discussion-conclusion}.

\end{document}
